\subsubsection{Graph tool : un graphe de dépendances structurelles}

Le \texttt{graph\_tool} complète le \textit{RAG} sémantique et la recherche lexicale par une troisième vue du code : une vue \textbf{structurelle}. L'idée n'est plus seulement de retrouver des morceaux de texte pertinents, mais de reconstruire explicitement les dépendances du projet sous la forme d'un graphe navigable.

À la différence du graphe d'orchestration \textit{LangGraph}, qui pilote la boucle de raisonnement de l'agent, le \texttt{graph\_tool} représente le contenu même du dépôt Envision. Chaque script, chaque dossier logique, chaque fichier de données, ainsi que certaines entités internes (tables, fonctions) deviennent des nœuds ; les relations \texttt{imports}, \texttt{reads}, \texttt{writes}, \texttt{defines}, \texttt{contains} et \texttt{sibling} deviennent des arêtes. Ce graphe permet donc de répondre à des questions structurelles du type : \textit{« quels scripts lisent \code{/Clean/Items.ion} ? »}, \textit{« quel module est importé par tel inspecteur ? »}, ou encore \textit{« dans quel script une table est-elle définie ? »}.

\paragraph{Deux arbres fondamentaux reliés}
Le graphe n'est pas un arbre unique. Il juxtapose deux hiérarchies distinctes :
\begin{itemize}
    \item le domaine \textbf{scripts}, qui organise les fichiers \texttt{.nvn} selon leur arborescence logique ;
    \item le domaine \textbf{data}, qui organise les fichiers \texttt{.ion}, \texttt{.csv} ou \texttt{.xlsx} manipulés par les scripts.
\end{itemize}
Ces deux arbres sont ensuite reliés par des arêtes transverses de lecture et d'écriture. Le graphe encode ainsi en une seule structure à la fois la hiérarchie du code et le flot des données.

\begin{figure}[H]
\begin{center}
\resizebox{0.96\textwidth}{!}{
\begin{tikzpicture}[ 
    >=Latex, 
    script/.style={ 
        rectangle, 
        rounded corners, 
        draw=blue!60!black, 
        fill=blue!10, 
        minimum height=0.95cm, 
        text width=4.4cm, 
        align=center, 
        inner sep=5pt 
    }, 
    data/.style={ 
        rectangle, 
        rounded corners, 
        draw=green!50!black, 
        fill=green!10, 
        minimum height=0.95cm, 
        text width=4.2cm, 
        align=center, 
        inner sep=5pt 
    }, 
    line/.style={draw, thick, -Latex}, 
    % Styles mis à jour : "text=" assure que le texte prend la même couleur
    relImport/.style={draw=orange!85!black, text=orange!85!black, thick, -Latex},
    relRead/.style={draw=red!70!black, text=red!70!black, dashed, thick, -Latex}
] 
    % --- NOEUDS SCRIPTS ---
    \node[script] (sroot) at (0,4.7) {\textbf{scripts : /}}; 
    
    \node[script] (s1) at (-5.2,3.0) {\code{/1. utilities}}; 
    \node[script] (s3) at (0,3.0) {\code{/3. Inspectors}}; 
    \node[script] (s4) at (5.2,3.0) {\shortstack{\code{/4. Optimization}\\\code{workflow}}}; 
    
    \node[script] (s1_mod) at (-5.2,1.3) {\shortstack{\code{/1. utilities/Modules/}\\\code{Global Parameters}}}; 
    \node[script] (s3_sales) at (0,1.3) {\shortstack{\code{/3. Inspectors/}\\\code{2 - Sales Analysis.nvn}}}; 

    % --- NOEUDS DATA ---
    \node[data] (droot) at (0,-1.0) {\textbf{data : /}}; 
    \node[data] (din) at (-5.2,-2.7) {\code{/Input}}; 
    \node[data] (dclean) at (0,-2.7) {\code{/Clean}}; 
    \node[data] (dman) at (5.2,-2.7) {\code{/Manual}}; 
    
    \node[data] (items) at (-5.2,-4.4) {\code{/Clean/Items.ion}}; 
    \node[data] (orders) at (0,-4.4) {\code{/Clean/Orders.ion}}; 
    \node[data] (stock) at (5.2,-4.4) {\code{/Clean/Stock.ion}}; 

    % --- LIENS STANDARDS (Noirs) ---
    \draw[line] (sroot) -- (s1); 
    \draw[line] (sroot) -- (s3); 
    \draw[line] (sroot) -- (s4); 
    
    \draw[line] (s1) -- (s1_mod);
    \draw[line] (s3) -- (s3_sales);

    \draw[line] (droot) -- (din); 
    \draw[line] (droot) -- (dclean); 
    \draw[line] (droot) -- (dman); 
    
    \draw[line] (dclean) -- (items); 
    \draw[line] (dclean) -- (orders); 
    \draw[line] (dclean) -- (stock); 

    % --- LIENS DE RELATIONS (Couleurs) ---
    % Imports (Orange)
    \draw[relImport] (s3_sales.west) -- node[above] {imports} (s1_mod.east); 
    
    % Reads (Rouges pointillés)
    \draw[relRead] (s3_sales.south west) -- node[left] {reads} (items.north); 
    \draw[relRead] (s3_sales.south east) -- node[right] {reads} (stock.north); 
    
    % Contournement fluide par la droite pour la flèche du milieu
    \draw[relRead] (s3_sales.south) .. controls +(2.6, -1.5) and +(2.6, 1.5) .. node[right] {reads} (orders.north); 
\end{tikzpicture}}
\end{center}
\caption{Visualisation conceptuelle du graphe interne : deux arbres fondamentaux (\textit{scripts} et \textit{data}) reliés par des arêtes de dépendance.}
\end{figure}

\paragraph{Exemple local de sous-graphe}
Une fois le graphe construit, l'agent peut raisonner localement autour d'un nœud particulier. La figure suivante montre le sous-graphe centré sur \code{/3. Inspectors/2 - Sales Analysis.nvn}, qui importe un module, lit plusieurs fichiers de données et définit une table.

\begin{figure}[H]
\begin{center}
\resizebox{0.92\textwidth}{!}{
\begin{tikzpicture}[
    >=Latex,
    node distance=1.6cm and 2.1cm,
    center/.style={
        shape=ellipse,
        draw=cyan!70!black,
        fill=cyan!18,
        minimum width=5.6cm,
        minimum height=1.55cm,
        text width=5.2cm,
        align=center,
        inner sep=6pt
    },
    script/.style={
        rectangle,
        rounded corners,
        draw=blue!60!black,
        fill=blue!10,
        minimum height=1cm,
        text width=4cm,
        align=center,
        inner sep=5pt
    },
    data/.style={
        rectangle,
        rounded corners,
        draw=green!50!black,
        fill=green!10,
        minimum height=1cm,
        text width=3.2cm,
        align=center,
        inner sep=5pt
    },
    internal/.style={
        rectangle,
        rounded corners,
        draw=orange!70!black,
        fill=orange!12,
        minimum height=1cm,
        text width=4cm,
        align=center,
        inner sep=5pt
    },
    line/.style={draw, thick, -Latex}
]
    \node[center] (sales) {\code{/3. Inspectors/2} \\ \code{- Sales Analysis.nvn}};
    \node[script, above left=1.8cm and 2.5cm of sales] (module) {\code{/1. utilities/Modules/} \\ {Global Parameters}};
    \node[data, right=3.1cm of sales, yshift=2.0cm] (items) {\code{/Clean/Items.ion}};
    \node[data, right=3.1cm of sales, yshift=0.7cm] (orders) {\code{/Clean/Orders.ion}};
    \node[data, right=3.1cm of sales, yshift=-0.7cm] (stock) {\code{/Clean/Stock.ion}};
    \node[data, right=3.1cm of sales, yshift=-2.0cm] (po) {\code{/Clean/PO.ion}};
    \node[internal, below left=1.8cm and 2.5cm of sales] (table) {\code{...::table::ItemsWeek}};

    \draw[line] (sales) -- node[pos=0.52, above, xshift=0.3cm]{\texttt{imports}} (module);
    \draw[line] (sales.east) -- node[pos=0.58, above, yshift=0.2cm]{\texttt{reads}} (items.west);
    \draw[line] (sales.east) -- node[pos=0.58, above]{\texttt{reads}} (orders.west);
    \draw[line] (sales.east) -- node[pos=0.58, below]{\texttt{reads}} (stock.west);
    \draw[line] (sales.east) -- node[pos=0.58, below, yshift=-0.2cm]{\texttt{reads}} (po.west);
    \draw[line] (sales) -- node[pos=0.52, below, xshift=0.4cm]{\texttt{defines}} (table);
\end{tikzpicture}}
\end{center}
\caption{Sous-graphe typique autour d'un script : le graphe fournit directement les relations structurelles utiles à l'agent.}
\end{figure}

\paragraph{Statistiques réelles du graphe}
Sur l'état courant du dépôt, le graphe généré le 11 avril 2026 contient 300 nœuds et 781 arêtes. La répartition est la suivante : 64 scripts, 73 fichiers de données, 96 tables, 17 fonctions et 50 dossiers. Côté relations, on compte notamment 55 arêtes \texttt{imports}, 270 arêtes \texttt{reads}, 63 arêtes \texttt{writes}, 113 arêtes \texttt{defines} et 185 arêtes \texttt{contains}. Cette densité montre bien que le \textit{graph tool} n'est pas un simple affichage d'arborescence, mais une représentation réellement relationnelle du dépôt.

\paragraph{Usage pratique}
Le graphe peut être interrogé indépendamment de l'agent via le CLI \code{python -m env_graph.network}. En pratique, ce CLI expose les mêmes familles d'opérations que l'outil agentique : lecture des statistiques globales, exploration d'un arbre \texttt{scripts} ou \texttt{data}, recherche d'un nœud par son nom, inspection de ses métadonnées, lecture locale de son contenu et navigation dans ses voisins entrants ou sortants. Il permet donc aussi bien une exploration manuelle du dépôt qu'une validation rapide de ce que l'agent utilise ensuite dans la boucle de raisonnement.

\paragraph{Apport pour l'agent}
Le \texttt{graph\_tool} apporte donc une forme de \textbf{mémoire structurelle précompilée}. Là où le \textit{RAG} sémantique répond bien aux questions de logique métier et où le \texttt{grep} répond bien aux recherches exactes, le graphe fournit une vérité relationnelle sur le dépôt :
\begin{itemize}
    \item qui lit quoi ;
    \item qui écrit quoi ;
    \item qui importe quoi ;
    \item dans quel dossier logique se trouve chaque élément ;
    \item quelles tables ou fonctions sont définies par quel script.
\end{itemize}

Cette couche est particulièrement précieuse dans Envision car de nombreux chemins sont construits dynamiquement par interpolation de chaînes. Le graphe capture et normalise une partie de ces dépendances en amont, ce qui réduit la charge d'inférence précédemment évoquée avec le \texttt{grep\_tool}.